%-------------------------------------------------------------------------------
\chapter{Dependency of the multijet background on $\Delta m_W$}
\label{a:mj_mW_dependency}
%--------------------------------

In general, fit results yield $\Delta m_W \neq 0$.
The multijet estimation is consistently performed using the data-driven Ansatz
%
\begin{equation}
    \mathrm{MJ} = \mathrm{Data} - \mathrm{EW}_{\mathrm{sim}},
    \label{eq:mj_ansatz}
\end{equation}
%
in which the electroweak (EW) signal simulation is generated under the assumption $\Delta m_W = 0$.
This assumption is inconsistent with the data, which by construction reflects the true, non-zero value of $\Delta m_W$.
Consequently, the resulting multijet template may retain residual $m_W$-dependent contributions, as the subtraction in Eq.~\ref{eq:mj_ansatz} relies on a signal simulation that does not incorporate the true mass value.

To quantify the magnitude of this effect, the following study is performed:
the multijet estimation is repeated using a signal template with $\Delta m_W$ shifted by $200\,\mathrm{MeV}$ relative to the nominal value.
The resulting multijet distributions and subsequent fit results are compared to those obtained with the nominal ($\Delta m_W = 0$) template.
This comparison provides a quantitative estimate of the sensitivity of the multijet background estimate to the mass hypothesis assumed in the subtraction procedure.
The study is performed using the $W \to e\nu$ channels, which are selected due to their large multijet background fraction and correspondingly enhanced sensitivity to this potential bias.
Figure~\ref{fig:mj_distributions} shows examples of comparisons between the obtained multijet distributions, normalised to their respective estimated yield.
The grey uncertainty bands associated with the shape and yield uncertainty cover the differences between the two estimations, hence no significant bias is observed.
%
\begin{figure}[htbp!]
    \centering
    \subfigure[]{\includegraphics[width=0.45\textwidth]{appendix/ptl_uT1_lEta1_cut8;1.pdf}}
    \subfigure[]{\includegraphics[width=0.45\textwidth]{appendix/ptl_uT4_lEta3_cut8;1.pdf}}
    
    \subfigure[]{\includegraphics[width=0.45\textwidth]{appendix/mT_uT1_lEta1_cut8;1.pdf}} 
    \subfigure[]{\includegraphics[width=0.45\textwidth]{appendix/mT_uT4_lEta3_cut8;1.pdf}}
    \caption{Comparison of MJ estimations for different assumed valus of $\Delta m_W$.
    The grey uncertainty band is calculated as the quadratic sum of the yield and shape uncertainty.
    The uncertainties cover the differences of the distributions.}
    \label{fig:mj_distributions}
\end{figure}
%
To quantify this behaviour further the impact of the assumed value in the multijet estimation is quantised with a fit study.
Figure~\ref{fig:mW_fits_mj} shows the fit results for $\Delta m_W$ for the two multijet estimations.
The impact of the different assumed values for $\Delta m_W$ in the multijet estimation on the extracted value for $\Delta m_W$ is small.
%
\begin{figure}[htbp!]
    \centering
    \subfigure[]{\includegraphics[width=0.45\textwidth]{appendix/0_CH1Result_Nominal_ptl.pdf}}
    \subfigure[]{\includegraphics[width=0.45\textwidth]{appendix/0_CH5Result_Nominal_ptl.pdf}}
    
    \subfigure[]{\includegraphics[width=0.45\textwidth]{appendix/200_CH1Result_Nominal_ptl.pdf}}
    \subfigure[]{\includegraphics[width=0.45\textwidth]{appendix/200_CH5Result_Nominal_ptl.pdf}}
    \caption{Fit results for $\Delta m_W$ using \ptl distributions.
    The upper row shows the result using a multijet estimation that assumes $\Delta m_W =0$, while the lower row corresponds to a multijet estimation using $\Delta m_W = 200 \, MeV$.
    The left column shows results for the $5\,TeV$ electron channel, the right column for the $13\, TeV$ electron channel.
    The impact of the assumed value for $\Delta m_W$ in the multijet estimation is less than $1\,MeV$ on the final result.}
    \label{fig:mW_fits_mj}
\end{figure}
%
This study supports the conclusion that the $m_W = 0$ assumption used in the nominal multijet estimation does not introduce a significant bias in the extracted $\Delta m_W$ value, assuming an initial shift of less than $200 \, MeV$.


